\section{Conclusion}
\label{sec:conclusion}

We presented \textbf{FC-GNN}, the first method to simultaneously provide
(i) fuzzy-rule-based GNN aggregation for uncertainty-aware cybersecurity
anomaly detection, (ii) distribution-free, community-conditional coverage
guarantees via Fuzzy Mondrian Conformal Prediction, and (iii) human-interpretable
IF-THEN rule explanations aligned with MITRE ATT\&CK.
Experiments on seven cybersecurity benchmarks confirm that FC-GNN is the
only evaluated method to reliably maintain $\ge\!90\%$ empirical coverage
with conditional validity across graph communities, while achieving prediction-set
efficiency competitive with or superior to the leading Mondrian CP baseline
(RR-GNN) on three datasets.
A formal coverage theorem establishes the finite-sample conditional
coverage slack as $O(|Q_m|^{-1/2})$ under the fuzzy permutation-invariance
condition.

\paragraph{Future work.}
Three directions are most promising:
(1) \emph{Federated FC-GNN} — extending FMCP to the federated GNN setting
\citep{akgul25,dang24federated} where calibration data is distributed
across clients without centralisation;
(2) \emph{Temporal FC-GNN} — incorporating non-exchangeable CP
\citep{ncpnet25} to provide valid coverage under concept drift without
retraining from scratch;
(3) \emph{Per-community Sugeno calibration} — learning community-specific
$\lambda$ parameters to tighten the coverage gap towards zero in all
communities simultaneously.

\begin{acknowledgements}
The author thanks the open-source maintainers of PyTorch Geometric, NetworkX,
and the python-louvain community detection library.
\end{acknowledgements}
